% =============================================================================
%  MSc Thesis — CI-Lib: A Computational Intelligence Library
% =============================================================================
\documentclass[12pt,a4paper,twoside]{report}

% ---- Packages ---------------------------------------------------------------
\usepackage[utf8]{inputenc}
\usepackage[T1]{fontenc}
\usepackage{geometry}
\geometry{margin=1in}
\usepackage{setspace}
\onehalfspacing
\usepackage{graphicx}
\usepackage{booktabs}
\usepackage{amsmath,amssymb}
\usepackage{listings}
\usepackage{xcolor}
\usepackage{natbib}
\usepackage{hyperref}
\usepackage{caption}
\usepackage{subcaption}
\usepackage{titlesec}
\usepackage{fancyhdr}
\usepackage{algorithm}
\usepackage{algpseudocode}

% ---- Code listing style ----------------------------------------------------
\lstset{
  language=Python,
  basicstyle=\small\ttfamily,
  keywordstyle=\color{blue},
  stringstyle=\color{red},
  commentstyle=\color{green!40!black},
  numbers=left,
  numberstyle=\tiny,
  frame=single,
  breaklines=true,
  showstringspaces=false,
}

% ---- Header / Footer -------------------------------------------------------
\pagestyle{fancy}
\fancyhf{}
\fancyhead[LE,RO]{\thepage}
\fancyhead[RE]{\leftmark}
\fancyhead[LO]{\rightmark}

% ---- Title formatting ------------------------------------------------------
\titleformat{\chapter}{\normalfont\huge\bfseries}{\thechapter.}{20pt}{\huge}

% =============================================================================
\begin{document}

% ---- Title page ------------------------------------------------------------
\begin{titlepage}
  \begin{center}
    \vspace*{2cm}
    {\Huge \bfseries CI-Lib: A Computational Intelligence Library}\\[0.5cm]
    {\Large A Pure-NumPy Framework for Optimisation, Learning,\\and Decision-Making}\\[2cm]
    {\large A Thesis submitted in partial fulfilment\\of the requirements for the degree of}\\[0.5cm]
    {\Large \bfseries Master of Science in Computer Science}\\[2cm]
    {\large Author Name}\\[0.3cm]
    {\large Student ID}\\[1cm]
    {\large Supervisor Name}\\[0.3cm]
    {\large Department of Computer Science}\\[0.3cm]
    {\large University Name}\\[2cm]
    {\large \today}
  \end{center}
\end{titlepage}

% ---- Declaration -----------------------------------------------------------
\chapter*{Declaration}
\addcontentsline{toc}{chapter}{Declaration}
I hereby declare that this thesis is my own work and that, to the best of my
knowledge and belief, it contains no material previously published or written
by another person, except where due reference is made.

\vspace{1cm}
\hfill Author Name \hfill

% ---- Abstract --------------------------------------------------------------
\chapter*{Abstract}
\addcontentsline{toc}{chapter}{Abstract}
This thesis presents CI-Lib, a pure-NumPy computational intelligence library
that implements seven core domains: neural networks, evolutionary algorithms,
swarm intelligence, fuzzy logic systems, clustering, classical optimisation,
and utility functions for preprocessing, evaluation, and benchmarking.

The library is designed with three design principles:
(1) zero dependencies beyond NumPy for the core algorithms,
(2) a consistent object-oriented API across all modules, and
(3) layered interfaces including a FastAPI REST backend, a Streamlit
interactive dashboard, and a cybersecurity incident response CLI (Tank \& Dozer).

Experimental evaluations demonstrate that each algorithm converges to known
optima on standard benchmark functions, validates correctness on canonical
problems (XOR for neural networks, the tipping problem for fuzzy inference,
the travelling salesman problem for ant colony optimisation), and achieves
competitive performance against reference implementations.

% ---- Acknowledgements ------------------------------------------------------
\chapter*{Acknowledgements}
\addcontentsline{toc}{chapter}{Acknowledgements}
To be completed.

% ---- Table of contents -----------------------------------------------------
\tableofcontents
\listoffigures
\listoftables

% =============================================================================
%  Chapters
% =============================================================================

\chapter{Introduction}
\label{chap:introduction}

\section{Motivation}
Computational Intelligence (CI) encompasses a family of nature-inspired
algorithms for solving complex optimisation, learning, and decision-making
problems \citep{engelbrecht2007computational}. While numerous CI libraries
exist (e.g., scikit-learn \citep{scikit-learn}, DEAP \citep{deap}, PyTorch
\citep{pytorch}), they often carry heavy dependencies, require specific
runtime environments, or abstract implementation details behind high-level
interfaces.

CI-Lib addresses these limitations by providing a clean, minimal-dependency,
pedagogically transparent implementation of core CI algorithms. The library
is intended for researchers, educators, and practitioners who need to
understand, modify, or extend CI algorithms without the overhead of large
frameworks.

\section{Contributions}
The principal contributions of this thesis are:
\begin{enumerate}
  \item A pure-NumPy computational intelligence library with seven modules
        covering neural networks, evolutionary algorithms, swarm intelligence,
        fuzzy logic, clustering, optimisation, and utilities.
  \item A layered architecture that separates algorithmic implementation
        from interface concerns, providing a FastAPI REST API, a Streamlit
        dashboard, and a CLI application.
  \item Empirical validation of all algorithms on standard benchmark functions
        and canonical test problems.
  \item Tank \& Dozer, a demonstration CLI framework that applies CI algorithms
        to cybersecurity incident response tasks.
\end{enumerate}

\section{Thesis Outline}
The remainder of this thesis is organised as follows:
Chapter~\ref{chap:background} reviews relevant literature;
Chapter~\ref{chap:design} describes the system architecture and design;
Chapter~\ref{chap:implementation} details the implementation of each module;
Chapter~\ref{chap:experiments} presents experimental results;
Chapter~\ref{chap:conclusion} concludes and discusses future work.


\chapter{Literature Review}
\label{chap:background}

\section{Computational Intelligence}
Computational Intelligence is a sub-field of artificial intelligence that
draws inspiration from biological and natural systems
\citep{engelbrecht2007computational}. It encompasses several paradigms:
artificial neural networks, evolutionary computation, swarm intelligence,
and fuzzy systems.

\section{Neural Networks}
Feedforward neural networks, trained via backpropagation
\citep{rumelhart1986learning}, form the foundation of modern deep learning.
The universal approximation theorem \citep{hornik1989multilayer} guarantees
that a feedforward network with a single hidden layer can approximate any
continuous function given sufficient neurons.

\section{Evolutionary Algorithms}
Genetic algorithms \citep{holland1975adaptation} and differential evolution
\citep{storn1997differential} are population-based optimisation methods that
mimic natural selection. They are particularly effective for non-convex,
multimodal, and discontinuous objective functions.

\section{Swarm Intelligence}
Particle swarm optimisation \citep{kennedy1995pso} and ant colony optimisation
\citep{dorigo1996antsystem} are decentralised, self-organising algorithms
inspired by social behaviour in nature.

\section{Fuzzy Logic}
Fuzzy set theory \citep{zadeh1965fuzzy} provides a mathematical framework
for representing uncertainty and imprecision. Mamdani fuzzy inference systems
\citep{mamdani1975experiment} are widely used in control and decision-making.

\section{Clustering}
K-Means \citep{lloyd1982least} and DBSCAN \citep{ester1996dbscan} are
fundamental unsupervised learning algorithms for discovering structure in
unlabelled data.

\section{Classical Optimisation}
Simulated annealing \citep{kirkpatrick1983optimization} and gradient-based
methods \citep{kingma2015adam} represent complementary approaches to
continuous optimisation: global exploration via stochastic search and local
refinement via gradient descent.


\chapter{System Architecture}
\label{chap:design}

\section{Design Principles}
CI-Lib follows three core design principles:
\begin{enumerate}
  \item \textbf{Minimal dependencies:} The core library depends only on NumPy,
        ensuring portability and ease of installation.
  \item \textbf{Consistent API:} All algorithms follow a uniform pattern:
        constructor with hyperparameters, a main computation method
        (e.g., \texttt{fit}, \texttt{evolve}, \texttt{optimize}), and
        access to results via public attributes.
  \item \textbf{Layered architecture:} The core library is separated from
        interface layers, allowing independent development and testing.
\end{enumerate}

\section{Architecture Overview}
Figure~\ref{fig:architecture} shows the system architecture.

\begin{figure}[h]
  \centering
  \includegraphics[width=0.8\textwidth]{figures/architecture.png}
  \caption{CI-Lib system architecture showing the three interface layers
           (REST API, Streamlit dashboard, Tank \& Dozer CLI) built on
           the core ci\_lib algorithms.}
  \label{fig:architecture}
\end{figure}

\section{Module Structure}
The library is organised into seven modules:

\begin{table}[h]
  \centering
  \begin{tabular}{lll}
    \toprule
    Module & Algorithms & Key Classes \\
    \midrule
    \texttt{neural} & Feedforward networks & \texttt{FeedForwardNetwork} \\
    \texttt{evolutionary} & GA, DE & \texttt{GeneticAlgorithm, DifferentialEvolution} \\
    \texttt{swarm} & PSO, ACO & \texttt{ParticleSwarmOptimizer, AntColonyOptimizer} \\
    \texttt{fuzzy} & Mamdani FIS & \texttt{FuzzySet, FuzzyVariable, FuzzyRule, FIS} \\
    \texttt{clustering} & K-Means, DBSCAN & \texttt{KMeans, DBSCAN} \\
    \texttt{optimization} & SA, GD & \texttt{SimulatedAnnealing, GradientDescent} \\
    \texttt{utils} & Preprocessing, metrics, benchmarks & \texttt{normalize, accuracy, sphere, ...} \\
    \bottomrule
  \end{tabular}
  \caption{CI-Lib module summary.}
  \label{tab:modules}
\end{table}


\chapter{Implementation}
\label{chap:implementation}

\section{Neural Networks}
The \texttt{FeedForwardNetwork} class implements a fully-connected network
with configurable layer sizes and activation functions. Weight initialisation
uses Xavier uniform initialisation \citep{glorot2010difficulty} to maintain
variance across layers.

Training uses mini-batch stochastic gradient descent with backpropagation.
The implementation supports sigmoid, tanh, ReLU, leaky ReLU, and softmax
activations.

\begin{algorithm}[h]
  \caption{Feedforward Network Training}\label{alg:nn}
  \begin{algorithmic}[1]
    \State Initialise weights $W^{(l)}$, biases $b^{(l)}$ via Xavier init
    \For{epoch $= 1$ to $E$}
      \For{each mini-batch $(X_b, y_b)$}
        \State Forward pass: $a^{(0)} = X_b$, $a^{(l)} = \sigma(z^{(l)})$,
              $z^{(l)} = a^{(l-1)}W^{(l)} + b^{(l)}$
        \State Compute loss: $\mathcal{L} = \frac{1}{m}\sum\|a^{(L)} - y_b\|^2$
        \State Backward pass: $\delta^{(L)} = a^{(L)} - y_b$,
              $\delta^{(l)} = (\delta^{(l+1)}{W^{(l+1)}}^\top) \odot \sigma'(z^{(l)})$
        \State Update: $W^{(l)} \gets W^{(l)} - \eta \cdot {a^{(l-1)}}^\top \delta^{(l)}$,
              $b^{(l)} \gets b^{(l)} - \eta \cdot \sum \delta^{(l)}$
      \EndFor
    \EndFor
  \end{algorithmic}
\end{algorithm}

\section{Evolutionary Algorithms}
Both the genetic algorithm and differential evolution follow the generational
loop pattern: initialise population, evaluate fitness, apply variation
operators (selection, crossover, mutation), and repeat.

The GA uses tournament selection, simulated binary crossover (SBX)
\citep{deb1995sbx}, Gaussian mutation, and elitism. The DE supports both
\texttt{best/1/bin} and \texttt{rand/1/bin} strategies with binomial crossover.

\section{Swarm Intelligence}
The PSO implementation uses inertia-weighted velocity updates with optional
linear decay \citep{shi1998modified}. The ACO implementation solves the
symmetric TSP using pheromone-based probabilistic route construction and
global pheromone update \citep{dorigo1996antsystem}.

\section{Fuzzy Logic}
The fuzzy logic module implements Mamdani-style inference with three
defuzzification methods: centroid, bisector, and mean-of-maxima.
Fuzzy sets support triangular, trapezoidal, and Gaussian membership
functions. Fuzzy operators (AND, OR, NOT) are overridden via Python's
dunder methods (\texttt{\&}, \texttt{|}, \texttt{\textasciitilde}).

\section{Clustering}
K-Means supports k-means++ initialisation and random initialisation, with
convergence based on centroid movement tolerance. DBSCAN uses BFS expansion
from core points to build clusters, labelling remaining points as noise.

\section{Optimisation}
Simulated annealing supports exponential, linear, and logarithmic cooling
schedules. Gradient descent implements three optimisers: vanilla SGD,
momentum \citep{rumelhart1986learning}, and Adam \citep{kingma2015adam},
with optional numerical gradient approximation via central finite differences.

\section{Interface Layers}
Three interface layers are built on the core library:
\begin{itemize}
  \item \textbf{FastAPI REST API} (Chapter~\ref{chap:design}) exposes all
        algorithms via 16 endpoints at \texttt{/api/*}. The API auto-generates
        OpenAPI documentation at \texttt{/docs}.
  \item \textbf{Streamlit Dashboard} provides an interactive, tabbed web
        interface for visualising algorithm behaviour in real time.
  \item \textbf{Tank \& Dozer CLI} demonstrates a practical application of CI
        algorithms to cybersecurity incident response, including log anomaly
        detection, IOC clustering, and report generation.
\end{itemize}


\chapter{Experimental Evaluation}
\label{chap:experiments}

\section{Benchmark Functions}
The library includes six standard benchmark functions for optimisation
algorithms: Sphere, Rosenbrock, Rastrigin, Ackley, Griewank, and Schwefel.
Each has a known global minimum, allowing quantitative evaluation of
convergence accuracy.

\section{Convergence Results}
Figure~\ref{fig:convergence} shows representative convergence behaviour
across algorithms on the Sphere function in 10 dimensions.

\begin{figure}[h]
  \centering
  \includegraphics[width=0.75\textwidth]{figures/convergence_comparison.png}
  \caption{Convergence comparison of GA, DE, PSO, and SA on the 10-D Sphere
           function. All algorithms approach the known optimum of 0.}
  \label{fig:convergence}
\end{figure}

\begin{table}[h]
  \centering
  \begin{tabular}{lccc}
    \toprule
    Algorithm & Sphere (10-D) & Rastrigin (10-D) & Rosenbrock (10-D) \\
    \midrule
    GA & \texttt{<result>} & \texttt{<result>} & \texttt{<result>} \\
    DE & \texttt{<result>} & \texttt{<result>} & \texttt{<result>} \\
    PSO & \texttt{<result>} & \texttt{<result>} & \texttt{<result>} \\
    SA & \texttt{<result>} & \texttt{<result>} & \texttt{<result>} \\
    \bottomrule
  \end{tabular}
  \caption{Best fitness achieved by each algorithm after 500 generations/iterations.}
  \label{tab:results}
\end{table}

\section{Neural Network Validation}
The feedforward network successfully learns the XOR problem, reducing MSE
from approximately 0.25 to below 0.01 within 2000 epochs with a 2-8-1
architecture.

\section{Fuzzy Inference Validation}
The Mamdani fuzzy inference system produces interpretable, monotonic
outputs for the classic tipping problem: higher service and food ratings
yield higher suggested tip percentages.

\section{Clustering Validation}
K-Means correctly identifies three well-separated Gaussian clusters with
silhouette scores exceeding 0.7. DBSCAN correctly labels isolated noise
points while clustering dense regions.

\section{ACO on TSP}
The ant colony optimiser finds near-optimal tours for randomly generated
TSP instances with up to 30 cities, consistently improving over random
baseline tours.


\chapter{Conclusion and Future Work}
\label{chap:conclusion}

\section{Summary}
This thesis presented CI-Lib, a pure-NumPy computational intelligence
library spanning seven core domains. The library's minimal-dependency
design, consistent API, and layered architecture make it suitable for
research, education, and practical applications.

\section{Future Work}
Several directions for future development include:
\begin{itemize}
  \item GPU acceleration via CuPy or JAX for larger-scale experiments.
  \item Additional CI paradigms, including learning classifier systems,
        artificial immune systems, and deep reinforcement learning.
  \item Integration with real-time data sources for the Tank \& Dozer CLI.
  \item Hyperparameter optimisation and automated algorithm configuration.
  \item Containerised deployment with Kubernetes orchestration.
\end{itemize}

% =============================================================================
%  Bibliography
% =============================================================================
\bibliographystyle{plainnat}
\bibliography{references}

% =============================================================================
%  Appendices
% =============================================================================
\appendix
\chapter{API Reference}
\label{app:api}
Full API documentation is auto-generated and available at \texttt{/docs}
when running the FastAPI backend.

\chapter{Installation Guide}
\label{app:installation}
See the \texttt{README.md} file or run:
\begin{lstlisting}
pip install -e ".[all]"
python -m backend.app          # start API
streamlit run frontend/app.py  # start dashboard
\end{lstlisting}

\end{document}
